%%%%%%%%%%%%%%%%
%%% deut 33 %%%
%%%%%%%%%%%%%%%%

\Chapter{33}

\Verse{1}
{
	\hDtr1{וְזֹאת הַבְּרָכָה אֲשֶׁר בֵּרַךְ מֹשֶׁה אִישׁ הָאֱלֹהִים אֶת בְּנֵי יִשְׂרָאֵל לִפְנֵי מוֹתוֹ׃}
}
{
	\eDtr1{
		And this is the blessing with which Moses, the man of God,
		blessed the children of Israel before his death.
	}
}
{
}

\Verse{2}
{
	\hDtr1{וַיֹּאמַר יְהֹוָה מִסִּינַי בָּא וְזָרַח מִשֵּׂעִיר לָמוֹ הוֹפִיעַ מֵהַר פָּארָן וְאָתָה מֵרִבְבֹת קֹדֶשׁ מִימִינוֹ (אשדת) [אֵשׁ דָּת] לָמוֹ׃}
}
{
	\eDtr1{
		And he said: YHWH came from Sinai and rose from Seir for
		them. He shone from Mount Paran and came from ten
		thousands of the holy, slopes at His right, for them.
	}
}
{
%	\footnote{
%		slopes at His right. No one knows what this line means. (Frank Cross and
%		David Noel Freedman once wrote, “Conjectures are almost as numerous as
%		scholars.”) I followed the consonantal text, taking Hebrew ’[image]dt as
%		“slopes,” which is what it means in its only other occurrences in the
%		Torah (Deut 3:17; 4:49). In both of those occurrences it refers to the
%		slopes of Pisgah, east of the Dead Sea. In this verse, YHWH shines like
%		the sun coming from the east over Seir and Paran, in which case the
%		slopes of Pisgah would in fact be at His right.
%	}
}

\Verse{3}
{
	\hDtr1{אַף חֹבֵב עַמִּים כּל קְדֹשָׁיו בְּיָדֶךָ וְהֵם תֻּכּוּ לְרַגְלֶךָ יִשָּׂא מִדַּבְּרֹתֶיךָ׃}
}
{
	\eDtr1{
		Also loving peoples, all holy ones are in your hand. And
		they knelt at your feet; they bore your words.
	}
}
{
}

\Verse{4}
{
	\hDtr1{תּוֹרָה צִוָּה לָנוּ מֹשֶׁה מוֹרָשָׁה קְהִלַּת יַעֲקֹב׃}
}
{
	\eDtr1{
		Moses commanded us instruction, a possession, community of
		Jacob.
	}
}
{
}

\Verse{5}
{
	\hDtr1{וַיְהִי בִישֻׁרוּן מֶלֶךְ בְּהִתְאַסֵּף רָאשֵׁי עָם יַחַד שִׁבְטֵי יִשְׂרָאֵל׃}
}
{
	\eDtr1{
		And He was king in Jeshurun when the people’s heads were
		gathered: Israel’s tribes together.
	}
}
{
}

\Verse{6}
{
	\hDtr1{יְחִי רְאוּבֵן וְאַל יָמֹת וִיהִי מְתָיו מִסְפָּר׃}
}
{
	\eDtr1{
		Let Reuben live and not die, but his men will be few in
		number.
	}
}
{
%	\footnote{
%		few in number. The territory of Reuben was the most vulnerable of any
%		tribe: located east of the Dead Sea, cut off from the tribes that were
%		west of the Jordan. It virtually ceased to exist by the tenth century
%		B.C.E.
%	}
}

\Verse{7}
{
	\hDtr1{וְזֹאת לִיהוּדָה וַיֹּאמַר שְׁמַע יְהֹוָה קוֹל יְהוּדָה וְאֶל עַמּוֹ תְּבִיאֶנּוּ יָדָיו רָב לוֹ וְעֵזֶר מִצָּרָיו תִּהְיֶה׃}
}
{
	\eDtr1{
		And this for Judah—and he said: Hear, YHWH, Judah’s voice,
		and bring him to his people. With his hands he strove for
		himself, and you’ll be a help from its foes.
	}
}
{
}

\Verse{8}
{
	\hDtr1{וּלְלֵוִי אָמַר תֻּמֶּיךָ וְאוּרֶיךָ לְאִישׁ חֲסִידֶךָ אֲשֶׁר נִסִּיתוֹ בְּמַסָּה תְּרִיבֵהוּ עַל מֵי מְרִיבָה׃}
}
{
	\eDtr1{
		And for Levi he said: Your Tummim and your Urim are your
		faithful man’s, whom you tested at Massah, disputed with
		him over Meribah’s water.
	}
}
{
}

\Verse{9}
{
	\hDtr1{הָאֹמֵר לְאָבִיו וּלְאִמּוֹ לֹא רְאִיתִיו וְאֶת אֶחָיו לֹא הִכִּיר וְאֶת בָּנָו לֹא יָדָע כִּי שָׁמְרוּ אִמְרָתֶךָ וּבְרִיתְךָ יִנְצֹרוּ׃}
}
{
	\eDtr1{
		Who said of his father and his mother: “I haven’t seen
		him,” and didn’t recognize his brothers and didn’t know
		his children, for they observed what you said and kept
		your covenant.
	}
}
{
%	\footnote{
%		father … mother … brothers … children. This has been taken in a narrow
%		sense (by Rashi and others) to refer to the golden calf episode, in
%		which Moses tells the Levites to “kill, each man, his brother and, each
%		man, his neighbor and, each man, his relative” (Exod 32:27,29). And it
%		has been understood more broadly, to refer to the Levites’ single-minded
%		performance of their priestly duties, neither being distracted by one’s
%		family nor showing favoritism to them. Alternatively, it might refer to
%		the requirement that the high priest may not touch any dead bodies, even
%		of his own parents (Lev 21:11).
%	}
}

\Verse{10}
{
	\hDtr1{יוֹרוּ מִשְׁפָּטֶיךָ לְיַעֲקֹב וְתוֹרָתְךָ לְיִשְׂרָאֵל יָשִׂימוּ קְטוֹרָה בְּאַפֶּךָ וְכָלִיל עַל מִזְבְּחֶךָ׃}
}
{
	\eDtr1{
		They’ll teach your judgments to Jacob and your instruction
		to Israel. They’ll set incense at your nose, entirely
		burnt on your altar.
	}
}
{
%	\footnote{
%		entirely burnt on your altar. The term for “entirely burnt”
%		(k[image]lîl) is used in connection with the offering that the priests
%		make on the day of the high priest’s anointing (Lev 6:12–16). It thus
%		especially symbolizes the priestly status of Levi.
%	}
}

\Verse{11}
{
	\hDtr1{בָּרֵךְ יְהֹוָה חֵילוֹ וּפֹעַל יָדָיו תִּרְצֶה מְחַץ מתְנַיִם קָמָיו וּמְשַׂנְאָיו מִן יְקוּמוּן׃}
}
{
	\eDtr1{
		Bless, YHWH, his wealth and accept his hands’ work. Pierce
		his adversaries’ hips, and those who hate him, so they
		won’t get up.
	}
}
{
}

\Verse{12}
{
	\hDtr1{לְבִנְיָמִן אָמַר יְדִיד יְהֹוָה יִשְׁכֹּן לָבֶטַח עָלָיו חֹפֵף עָלָיו כּל הַיּוֹם וּבֵין כְּתֵפָיו שָׁכֵן׃}
}
{
	\eDtr1{
		For Benjamin he said: Beloved of YHWH, he’ll dwell in
		security by Him. He shelters over him all day as he dwells
		between his shoulders.
	}
}
{
}

\Verse{13}
{
	\hDtr1{וּלְיוֹסֵף אָמַר מְבֹרֶכֶת יְהֹוָה אַרְצוֹ מִמֶּגֶד שָׁמַיִם מִטָּל וּמִתְּהוֹם רֹבֶצֶת תָּחַת׃}
}
{
	\eDtr1{
		And for Joseph he said: Blessed of YHWH is his land, from
		the skies’ abundance, from dew, and from the deep,
		crouching below,
	}
}
{
}

\Verse{14}
{
	\hDtr1{וּמִמֶּגֶד תְּבוּאֹת שָׁמֶשׁ וּמִמֶּגֶד גֶּרֶשׁ יְרָחִים׃}
}
{
	\eDtr1{
		and from the abundance of the sun’s produce and from the
		abundance of the moon’s output
	}
}
{
}

\Verse{15}
{
	\hDtr1{וּמֵרֹאשׁ הַרְרֵי קֶדֶם וּמִמֶּגֶד גִּבְעוֹת עוֹלָם׃}
}
{
	\eDtr1{
		and from the top of the ancient mountains and from the
		abundance of the hills of antiquity
	}
}
{
}

\Verse{16}
{
	\hDtr1{וּמִמֶּגֶד אֶרֶץ וּמְלֹאָהּ וּרְצוֹן שֹׁכְנִי סְנֶה תָּבוֹאתָה לְרֹאשׁ יוֹסֵף וּלְקדְקֹד נְזִיר אֶחָיו׃}
}
{
	\eDtr1{
		and from the abundance of earth and what fills it and the
		favor of the one who dwelt in the bush. May it be on
		Joseph’s head, on the top of the head of the one separate
		from his brothers.
	}
}
{
%	\footnote{
%		May it be. Following Cross and Freedman, I have read the Hebrew [image]
%		on the likelihood that it conflated two readings: “May it come”
%		([image]) and “May it be” ([image]). The latter is preferable on the
%		basis of comparison with the parallel blessing of Joseph in Gen 49:26.
%	}
}

\Verse{17}
{
	\hDtr1{בְּכוֹר שׁוֹרוֹ הָדָר לוֹ וְקַרְנֵי רְאֵם קַרְנָיו בָּהֶם עַמִּים יְנַגַּח יַחְדָּו אַפְסֵי אָרֶץ וְהֵם רִבְבוֹת אֶפְרַיִם וְהֵם אַלְפֵי מְנַשֶּׁה׃}
}
{
	\eDtr1{
		His firstborn bull: it has majesty. And its horns are a
		wild ox’s horns. It will gore peoples with them, together,
		the ends of the earth. And they’re Ephraim’s ten
		thousands, And they’re Manasseh’s thousands.
	}
}
{
}

\Verse{18}
{
	\hDtr1{וְלִזְבוּלֻן אָמַר שְׂמַח זְבוּלֻן בְּצֵאתֶךָ וְיִשָּׂשכָר בְּאֹהָלֶיךָ׃}
}
{
	\eDtr1{
		And for Zebulun he said: Rejoice, Zebulun, when you go
		out, and Issachar in your tents.
	}
}
{
}

\Verse{19}
{
	\hDtr1{עַמִּים הַר יִקְרָאוּ שָׁם יִזְבְּחוּ זִבְחֵי צֶדֶק כִּי שֶׁפַע יַמִּים יִינָקוּ וּשְׂפֻנֵי טְמוּנֵי חוֹל׃}
}
{
	\eDtr1{
		They’ll call peoples to the mountain. There they’ll offer
		the sacrifices of virtue. For they’ll suck the seas’
		bounty and the sand’s hidden treasures.
	}
}
{
}

\Verse{20}
{
	\hDtr1{וּלְגָד אָמַר בָּרוּךְ מַרְחִיב גָּד כְּלָבִיא שָׁכֵן וְטָרַף זְרוֹעַ אַף קדְקֹד׃}
}
{
	\eDtr1{
		And for Gad he said: Blessed is one who enlarges Gad. Like
		a feline, abiding and tearing an arm and the top of a
		head,
	}
}
{
}

\Verse{21}
{
	\hDtr1{וַיַּרְא רֵאשִׁית לוֹ כִּי שָׁם חֶלְקַת מְחֹקֵק סָפוּן וַיֵּתֵא רָאשֵׁי עָם צִדְקַת יְהֹוָה עָשָׂה וּמִשְׁפָּטָיו עִם יִשְׂרָאֵל׃}
}
{
	\eDtr1{
		so he saw the foremost for himself, for a ruler’s share
		was kept there. And the heads of the people came. And he
		did YHWH’s justice and His laws with Israel.
	}
}
{
%	\footnote{
%		he saw the foremost for himself. Neither the length nor the content of
%		this verse suits the tribe of Gad, which is covered in the previous
%		verse. My colleague David Noel Freedman has suggested that it is rather
%		intertribal and that the specific references are to Moses as the leader
%		of the confederation of the tribes. This matches the picture of Moses in
%		vv. 4–5.
%	}
}

\Verse{22}
{
	\hDtr1{וּלְדָן אָמַר דָּן גּוּר אַרְיֵה יְזַנֵּק מִן הַבָּשָׁן׃}
}
{
	\eDtr1{
		And for Dan he said: Dan is a lion’s whelp that leapt from
		Bashan.
	}
}
{
%	\footnote{
%		Dan is a lion’s whelp. But in Jacob’s blessing “Judah is a lion’s
%		whelp”! (Gen 49:9). Why are these two tribes described this way? Dan is
%		the northernmost tribe, and Judah is the southernmost. So this produces
%		an image of Israel protected by lions on either side. (Actually, Simeon
%		is the southernmost tribe, but Simeon is unaccountably left out of this
%		song.)
%	}
}

\Verse{23}
{
	\hDtr1{וּלְנַפְתָּלִי אָמַר נַפְתָּלִי שְׂבַע רָצוֹן וּמָלֵא בִּרְכַּת יְהֹוָה יָם וְדָרוֹם יְרָשָׁה׃}
}
{
	\eDtr1{
		And for Naphtali he said: Naphtali is full of favor and
		filled with YHWH’s blessing. He’ll possess west and south.
	}
}
{
%	\footnote{
%		west and south. But Naphtali is located in the north and east of Israel.
%		So this verse is understood to refer to the west and south of the Sea of
%		Galilee (known also as Sea of Tiberias, or Sea of Kinneret).
%	}
}

\Verse{24}
{
	\hDtr1{וּלְאָשֵׁר אָמַר בָּרוּךְ מִבָּנִים אָשֵׁר יְהִי רְצוּי אֶחָיו וְטֹבֵל בַּשֶּׁמֶן רַגְלוֹ׃}
}
{
	\eDtr1{
		And for Asher he said: Blessed out of the sons is Asher.
		Let him be favored by his brothers and dipping his foot in
		oil.
	}
}
{
}

\Verse{25}
{
	\hDtr1{בַּרְזֶל וּנְחֹשֶׁת מִנְעָלֶךָ וּכְיָמֶיךָ דּבְאֶךָ׃}
}
{
	\eDtr1{
		Your lock is iron and bronze, and your strength is as much
		as your days.
	}
}
{
}

\Verse{26}
{
	\hDtr1{אֵין כָּאֵל יְשֻׁרוּן רֹכֵב שָׁמַיִם בְּעֶזְרֶךָ וּבְגַאֲוָתוֹ שְׁחָקִים׃}
}
{
	\eDtr1{
		There’s none like the God of Jeshurun, riding skies to
		help you and clouds in His majesty.
	}
}
{
}

\Verse{27}
{
	\hDtr1{מְעֹנָה אֱלֹהֵי קֶדֶם וּמִתַּחַת זְרֹעֹת עוֹלָם וַיְגָרֶשׁ מִפָּנֶיךָ אוֹיֵב וַיֹּאמֶר הַשְׁמֵד׃}
}
{
	\eDtr1{
		The ancient God is a refuge; and below: the arms of the
		eternal. And He drove out an enemy before you and said,
		“Destroy!”
	}
}
{
}

\Verse{28}
{
	\hDtr1{וַיִּשְׁכֹּן יִשְׂרָאֵל בֶּטַח בָּדָד עֵין יַעֲקֹב אֶל אֶרֶץ דָּגָן וְתִירוֹשׁ אַף שָׁמָיו יַעַרְפוּ טָל׃}
}
{
	\eDtr1{
		And Israel will dwell secure, Jacob dwells alone. To a
		land of grain and wine; and its skies drop dew.
	}
}
{
}

\Verse{29}
{
	\hDtr1{אַשְׁרֶיךָ יִשְׂרָאֵל מִי כָמוֹךָ עַם נוֹשַׁע בַּיהֹוָה מָגֵן עֶזְרֶךָ וַאֲשֶׁר חֶרֶב גַּאֲוָתֶךָ וְיִכָּחֲשׁוּ אֹיְבֶיךָ לָךְ וְאַתָּה עַל בָּמוֹתֵימוֹ תִדְרֹךְ׃}
}
{
	\eDtr1{
		Happy are you, Israel! Who is like you, a people saved by
		YHWH, your strong shield who is your majestic sword?! And
		your enemies will fawn to you. And you: you’ll step on
		their high places.
	}
}
{
%	\footnote{
%		Happy are you, Israel! These are the last words that the people hear
%		from Moses. Strange: The first half is understandable as his final
%		words, “Happy are you … a people saved by YHWH.” The second half is not
%		what we might have expected: God is their strong shield and sword! “And
%		your enemies will fawn to you. And you: you’ll step on their high
%		places.” Why does he end with the defeat of enemies? Apparently an
%		assurance that Israel does not have to fear enemies was important—and so
%		it has in fact turned out in the people’s history in the three millennia
%		following Moses. The presence of enemies has been a sad, constant fact
%		of life. But Moses’ last words are an assurance that Israel will survive
%		to fulfill the destiny that was the first promise to Abraham: to be a
%		source of blessing for all the earth’s families. Footnote 33:29. high
%		places. See the comment on Lev 26:30.
%	}
}
